% Options for packages loaded elsewhere
\PassOptionsToPackage{unicode}{hyperref}
\PassOptionsToPackage{hyphens}{url}
\documentclass[
]{article}
\usepackage{xcolor}
\usepackage[margin=1in]{geometry}
\usepackage{amsmath,amssymb}
\setcounter{secnumdepth}{-\maxdimen} % remove section numbering
\usepackage{iftex}
\ifPDFTeX
  \usepackage[T1]{fontenc}
  \usepackage[utf8]{inputenc}
  \usepackage{textcomp} % provide euro and other symbols
\else % if luatex or xetex
  \usepackage{unicode-math} % this also loads fontspec
  \defaultfontfeatures{Scale=MatchLowercase}
  \defaultfontfeatures[\rmfamily]{Ligatures=TeX,Scale=1}
\fi
\usepackage{lmodern}
\ifPDFTeX\else
  % xetex/luatex font selection
\fi
% Use upquote if available, for straight quotes in verbatim environments
\IfFileExists{upquote.sty}{\usepackage{upquote}}{}
\IfFileExists{microtype.sty}{% use microtype if available
  \usepackage[]{microtype}
  \UseMicrotypeSet[protrusion]{basicmath} % disable protrusion for tt fonts
}{}
\makeatletter
\@ifundefined{KOMAClassName}{% if non-KOMA class
  \IfFileExists{parskip.sty}{%
    \usepackage{parskip}
  }{% else
    \setlength{\parindent}{0pt}
    \setlength{\parskip}{6pt plus 2pt minus 1pt}}
}{% if KOMA class
  \KOMAoptions{parskip=half}}
\makeatother
\usepackage{color}
\usepackage{fancyvrb}
\newcommand{\VerbBar}{|}
\newcommand{\VERB}{\Verb[commandchars=\\\{\}]}
\DefineVerbatimEnvironment{Highlighting}{Verbatim}{commandchars=\\\{\}}
% Add ',fontsize=\small' for more characters per line
\usepackage{framed}
\definecolor{shadecolor}{RGB}{248,248,248}
\newenvironment{Shaded}{\begin{snugshade}}{\end{snugshade}}
\newcommand{\AlertTok}[1]{\textcolor[rgb]{0.94,0.16,0.16}{#1}}
\newcommand{\AnnotationTok}[1]{\textcolor[rgb]{0.56,0.35,0.01}{\textbf{\textit{#1}}}}
\newcommand{\AttributeTok}[1]{\textcolor[rgb]{0.13,0.29,0.53}{#1}}
\newcommand{\BaseNTok}[1]{\textcolor[rgb]{0.00,0.00,0.81}{#1}}
\newcommand{\BuiltInTok}[1]{#1}
\newcommand{\CharTok}[1]{\textcolor[rgb]{0.31,0.60,0.02}{#1}}
\newcommand{\CommentTok}[1]{\textcolor[rgb]{0.56,0.35,0.01}{\textit{#1}}}
\newcommand{\CommentVarTok}[1]{\textcolor[rgb]{0.56,0.35,0.01}{\textbf{\textit{#1}}}}
\newcommand{\ConstantTok}[1]{\textcolor[rgb]{0.56,0.35,0.01}{#1}}
\newcommand{\ControlFlowTok}[1]{\textcolor[rgb]{0.13,0.29,0.53}{\textbf{#1}}}
\newcommand{\DataTypeTok}[1]{\textcolor[rgb]{0.13,0.29,0.53}{#1}}
\newcommand{\DecValTok}[1]{\textcolor[rgb]{0.00,0.00,0.81}{#1}}
\newcommand{\DocumentationTok}[1]{\textcolor[rgb]{0.56,0.35,0.01}{\textbf{\textit{#1}}}}
\newcommand{\ErrorTok}[1]{\textcolor[rgb]{0.64,0.00,0.00}{\textbf{#1}}}
\newcommand{\ExtensionTok}[1]{#1}
\newcommand{\FloatTok}[1]{\textcolor[rgb]{0.00,0.00,0.81}{#1}}
\newcommand{\FunctionTok}[1]{\textcolor[rgb]{0.13,0.29,0.53}{\textbf{#1}}}
\newcommand{\ImportTok}[1]{#1}
\newcommand{\InformationTok}[1]{\textcolor[rgb]{0.56,0.35,0.01}{\textbf{\textit{#1}}}}
\newcommand{\KeywordTok}[1]{\textcolor[rgb]{0.13,0.29,0.53}{\textbf{#1}}}
\newcommand{\NormalTok}[1]{#1}
\newcommand{\OperatorTok}[1]{\textcolor[rgb]{0.81,0.36,0.00}{\textbf{#1}}}
\newcommand{\OtherTok}[1]{\textcolor[rgb]{0.56,0.35,0.01}{#1}}
\newcommand{\PreprocessorTok}[1]{\textcolor[rgb]{0.56,0.35,0.01}{\textit{#1}}}
\newcommand{\RegionMarkerTok}[1]{#1}
\newcommand{\SpecialCharTok}[1]{\textcolor[rgb]{0.81,0.36,0.00}{\textbf{#1}}}
\newcommand{\SpecialStringTok}[1]{\textcolor[rgb]{0.31,0.60,0.02}{#1}}
\newcommand{\StringTok}[1]{\textcolor[rgb]{0.31,0.60,0.02}{#1}}
\newcommand{\VariableTok}[1]{\textcolor[rgb]{0.00,0.00,0.00}{#1}}
\newcommand{\VerbatimStringTok}[1]{\textcolor[rgb]{0.31,0.60,0.02}{#1}}
\newcommand{\WarningTok}[1]{\textcolor[rgb]{0.56,0.35,0.01}{\textbf{\textit{#1}}}}
\usepackage{graphicx}
\makeatletter
\newsavebox\pandoc@box
\newcommand*\pandocbounded[1]{% scales image to fit in text height/width
  \sbox\pandoc@box{#1}%
  \Gscale@div\@tempa{\textheight}{\dimexpr\ht\pandoc@box+\dp\pandoc@box\relax}%
  \Gscale@div\@tempb{\linewidth}{\wd\pandoc@box}%
  \ifdim\@tempb\p@<\@tempa\p@\let\@tempa\@tempb\fi% select the smaller of both
  \ifdim\@tempa\p@<\p@\scalebox{\@tempa}{\usebox\pandoc@box}%
  \else\usebox{\pandoc@box}%
  \fi%
}
% Set default figure placement to htbp
\def\fps@figure{htbp}
\makeatother
\setlength{\emergencystretch}{3em} % prevent overfull lines
\providecommand{\tightlist}{%
  \setlength{\itemsep}{0pt}\setlength{\parskip}{0pt}}
\usepackage{bookmark}
\IfFileExists{xurl.sty}{\usepackage{xurl}}{} % add URL line breaks if available
\urlstyle{same}
\hypersetup{
  pdftitle={Diagnostic testing: Se{,} Sp{,} PPV{,} NPV{,} LR},
  hidelinks,
  pdfcreator={LaTeX via pandoc}}

\title{Diagnostic testing: Se, Sp, PPV, NPV, LR}
\author{}
\date{\vspace{-2.5em}}

\begin{document}
\maketitle

\textbf{Inference labs} use the five-step template: Hypothesis →
Visualise → Assumptions → Conduct → Conclude.

\subsection{Learning objectives}\label{learning-objectives}

\begin{itemize}
\tightlist
\item
  Compute sensitivity, specificity, PPV, NPV, and positive and negative
  likelihood ratios from a 2x2 table.
\item
  Convert pre-test probability to post-test probability with an LR.
\item
  Sketch a receiver-operating characteristic curve from a continuous
  test statistic.
\end{itemize}

\subsection{Prerequisites}\label{prerequisites}

Lab 2.2.

\subsection{Background}\label{background}

A diagnostic test has two operating characteristics intrinsic to the
test itself: \emph{sensitivity} is the probability that a diseased
person tests positive; \emph{specificity} is the probability that a
disease-free person tests negative. These quantities are properties of
the test. They do not change with prevalence.

Two other quantities are properties of the test \emph{and} the
population in which it is applied: \emph{positive predictive value} is
the probability of disease given a positive test; \emph{negative
predictive value} is the probability of no disease given a negative
test. These change with prevalence, sometimes dramatically.

Likelihood ratios unify the two pairs. \emph{LR+} is sens / (1 − spec);
\emph{LR−} is (1 − sens) / spec. They convert pre-test odds to post-test
odds by multiplication, which is the cleanest way to combine a test
result with prior information. An LR+ greater than 10 is a strong
positive; less than 0.1 is a strong negative; values near 1 are
uninformative.

\subsection{Setup}\label{setup}

\begin{Shaded}
\begin{Highlighting}[]
\FunctionTok{library}\NormalTok{(tidyverse)}
\FunctionTok{set.seed}\NormalTok{(}\DecValTok{42}\NormalTok{)}
\FunctionTok{theme\_set}\NormalTok{(}\FunctionTok{theme\_minimal}\NormalTok{(}\AttributeTok{base\_size =} \DecValTok{12}\NormalTok{))}
\end{Highlighting}
\end{Shaded}

\subsection{1. Hypothesis}\label{hypothesis}

Question of interest: how does a continuous biomarker behave as a
diagnostic test? We are not running an inferential test; we are
characterising a test's discrimination.

\subsection{2. Visualise}\label{visualise}

Simulate a biomarker that is higher in diseased cases than in
disease-free controls, with overlap.

\begin{Shaded}
\begin{Highlighting}[]
\NormalTok{prev }\OtherTok{\textless{}{-}} \FloatTok{0.2}
\NormalTok{pop }\OtherTok{\textless{}{-}} \FunctionTok{tibble}\NormalTok{(}
  \AttributeTok{id =} \FunctionTok{seq\_len}\NormalTok{(N),}
  \AttributeTok{disease =} \FunctionTok{rbinom}\NormalTok{(N, }\DecValTok{1}\NormalTok{, prev),}
  \AttributeTok{biomarker =} \FunctionTok{rnorm}\NormalTok{(N, }\AttributeTok{mean =} \FunctionTok{if\_else}\NormalTok{(disease }\SpecialCharTok{==} \DecValTok{1}\NormalTok{, }\DecValTok{7}\NormalTok{, }\DecValTok{5}\NormalTok{), }\AttributeTok{sd =} \DecValTok{1}\NormalTok{)}
\NormalTok{)}
\end{Highlighting}
\end{Shaded}

\begin{Shaded}
\begin{Highlighting}[]
\NormalTok{pop }\SpecialCharTok{|\textgreater{}}
  \FunctionTok{mutate}\NormalTok{(}\AttributeTok{status =} \FunctionTok{if\_else}\NormalTok{(disease }\SpecialCharTok{==} \DecValTok{1}\NormalTok{, }\StringTok{"disease"}\NormalTok{, }\StringTok{"no disease"}\NormalTok{)) }\SpecialCharTok{|\textgreater{}}
  \FunctionTok{ggplot}\NormalTok{(}\FunctionTok{aes}\NormalTok{(biomarker, }\AttributeTok{fill =}\NormalTok{ status)) }\SpecialCharTok{+}
  \FunctionTok{geom\_density}\NormalTok{(}\AttributeTok{alpha =} \FloatTok{0.5}\NormalTok{, }\AttributeTok{colour =} \ConstantTok{NA}\NormalTok{) }\SpecialCharTok{+}
  \FunctionTok{geom\_vline}\NormalTok{(}\AttributeTok{xintercept =} \DecValTok{6}\NormalTok{, }\AttributeTok{linetype =} \DecValTok{2}\NormalTok{) }\SpecialCharTok{+}
  \FunctionTok{labs}\NormalTok{(}\AttributeTok{x =} \StringTok{"Biomarker level"}\NormalTok{, }\AttributeTok{y =} \StringTok{"Density"}\NormalTok{, }\AttributeTok{fill =} \ConstantTok{NULL}\NormalTok{)}
\end{Highlighting}
\end{Shaded}

\subsection{3. Assumptions}\label{assumptions}

The gold standard for disease status is assumed perfect. The biomarker
is continuous and must be dichotomised at some cutoff to behave like a
positive/negative test. We choose 6 as the cutoff for illustration; in
practice, the cutoff is itself an outcome of the analysis.

\begin{Shaded}
\begin{Highlighting}[]
\NormalTok{pop }\OtherTok{\textless{}{-}}\NormalTok{ pop }\SpecialCharTok{|\textgreater{}} \FunctionTok{mutate}\NormalTok{(}\AttributeTok{test =} \FunctionTok{as.integer}\NormalTok{(biomarker }\SpecialCharTok{\textgreater{}}\NormalTok{ cutoff))}
\NormalTok{tab }\OtherTok{\textless{}{-}} \FunctionTok{table}\NormalTok{(}\AttributeTok{disease =}\NormalTok{ pop}\SpecialCharTok{$}\NormalTok{disease, }\AttributeTok{test =}\NormalTok{ pop}\SpecialCharTok{$}\NormalTok{test)}
\NormalTok{tab}
\end{Highlighting}
\end{Shaded}

\subsection{4. Conduct}\label{conduct}

\begin{Shaded}
\begin{Highlighting}[]
\NormalTok{FP }\OtherTok{\textless{}{-}}\NormalTok{ tab[}\StringTok{"0"}\NormalTok{, }\StringTok{"1"}\NormalTok{]; TN }\OtherTok{\textless{}{-}}\NormalTok{ tab[}\StringTok{"0"}\NormalTok{, }\StringTok{"0"}\NormalTok{]}

\NormalTok{sens }\OtherTok{\textless{}{-}}\NormalTok{ TP }\SpecialCharTok{/}\NormalTok{ (TP }\SpecialCharTok{+}\NormalTok{ FN)}
\NormalTok{spec }\OtherTok{\textless{}{-}}\NormalTok{ TN }\SpecialCharTok{/}\NormalTok{ (TN }\SpecialCharTok{+}\NormalTok{ FP)}
\NormalTok{ppv  }\OtherTok{\textless{}{-}}\NormalTok{ TP }\SpecialCharTok{/}\NormalTok{ (TP }\SpecialCharTok{+}\NormalTok{ FP)}
\NormalTok{npv  }\OtherTok{\textless{}{-}}\NormalTok{ TN }\SpecialCharTok{/}\NormalTok{ (TN }\SpecialCharTok{+}\NormalTok{ FN)}
\NormalTok{lrp  }\OtherTok{\textless{}{-}}\NormalTok{ sens }\SpecialCharTok{/}\NormalTok{ (}\DecValTok{1} \SpecialCharTok{{-}}\NormalTok{ spec)}
\NormalTok{lrn  }\OtherTok{\textless{}{-}}\NormalTok{ (}\DecValTok{1} \SpecialCharTok{{-}}\NormalTok{ sens) }\SpecialCharTok{/}\NormalTok{ spec}

\NormalTok{diag\_tbl }\OtherTok{\textless{}{-}} \FunctionTok{tibble}\NormalTok{(}
  \AttributeTok{quantity =} \FunctionTok{c}\NormalTok{(}\StringTok{"Sensitivity"}\NormalTok{, }\StringTok{"Specificity"}\NormalTok{,}
               \StringTok{"PPV"}\NormalTok{, }\StringTok{"NPV"}\NormalTok{, }\StringTok{"LR+"}\NormalTok{, }\StringTok{"LR{-}"}\NormalTok{),}
  \AttributeTok{value =} \FunctionTok{c}\NormalTok{(sens, spec, ppv, npv, lrp, lrn)}
\NormalTok{)}
\NormalTok{diag\_tbl}
\end{Highlighting}
\end{Shaded}

Convert pre-test odds to post-test odds with the LR.

\begin{Shaded}
\begin{Highlighting}[]
\NormalTok{pre\_odds }\OtherTok{\textless{}{-}}\NormalTok{ pre\_prob }\SpecialCharTok{/}\NormalTok{ (}\DecValTok{1} \SpecialCharTok{{-}}\NormalTok{ pre\_prob)}
\NormalTok{post\_odds\_pos }\OtherTok{\textless{}{-}}\NormalTok{ pre\_odds }\SpecialCharTok{*}\NormalTok{ lrp}
\NormalTok{post\_prob\_pos }\OtherTok{\textless{}{-}}\NormalTok{ post\_odds\_pos }\SpecialCharTok{/}\NormalTok{ (}\DecValTok{1} \SpecialCharTok{+}\NormalTok{ post\_odds\_pos)}
\NormalTok{post\_odds\_neg }\OtherTok{\textless{}{-}}\NormalTok{ pre\_odds }\SpecialCharTok{*}\NormalTok{ lrn}
\NormalTok{post\_prob\_neg }\OtherTok{\textless{}{-}}\NormalTok{ post\_odds\_neg }\SpecialCharTok{/}\NormalTok{ (}\DecValTok{1} \SpecialCharTok{+}\NormalTok{ post\_odds\_neg)}

\FunctionTok{tibble}\NormalTok{(}
\NormalTok{  pre\_prob,}
  \AttributeTok{post\_prob\_if\_positive =}\NormalTok{ post\_prob\_pos,}
  \AttributeTok{post\_prob\_if\_negative =}\NormalTok{ post\_prob\_neg}
\NormalTok{)}
\end{Highlighting}
\end{Shaded}

Sketch an ROC by sweeping the cutoff.

\begin{Shaded}
\begin{Highlighting}[]
\NormalTok{  cut }\OtherTok{=} \FunctionTok{seq}\NormalTok{(}\FunctionTok{min}\NormalTok{(pop}\SpecialCharTok{$}\NormalTok{biomarker), }\FunctionTok{max}\NormalTok{(pop}\SpecialCharTok{$}\NormalTok{biomarker), }\AttributeTok{length.out =} \DecValTok{200}\NormalTok{)}
\ErrorTok{)} \SpecialCharTok{|\textgreater{}}
  \FunctionTok{rowwise}\NormalTok{() }\SpecialCharTok{|\textgreater{}}
  \FunctionTok{mutate}\NormalTok{(}
    \AttributeTok{tp =} \FunctionTok{sum}\NormalTok{(pop}\SpecialCharTok{$}\NormalTok{biomarker }\SpecialCharTok{\textgreater{}}\NormalTok{ cut }\SpecialCharTok{\&}\NormalTok{ pop}\SpecialCharTok{$}\NormalTok{disease }\SpecialCharTok{==} \DecValTok{1}\NormalTok{),}
    \AttributeTok{fn =} \FunctionTok{sum}\NormalTok{(pop}\SpecialCharTok{$}\NormalTok{biomarker }\SpecialCharTok{\textless{}=}\NormalTok{ cut }\SpecialCharTok{\&}\NormalTok{ pop}\SpecialCharTok{$}\NormalTok{disease }\SpecialCharTok{==} \DecValTok{1}\NormalTok{),}
    \AttributeTok{fp =} \FunctionTok{sum}\NormalTok{(pop}\SpecialCharTok{$}\NormalTok{biomarker }\SpecialCharTok{\textgreater{}}\NormalTok{ cut }\SpecialCharTok{\&}\NormalTok{ pop}\SpecialCharTok{$}\NormalTok{disease }\SpecialCharTok{==} \DecValTok{0}\NormalTok{),}
    \AttributeTok{tn =} \FunctionTok{sum}\NormalTok{(pop}\SpecialCharTok{$}\NormalTok{biomarker }\SpecialCharTok{\textless{}=}\NormalTok{ cut }\SpecialCharTok{\&}\NormalTok{ pop}\SpecialCharTok{$}\NormalTok{disease }\SpecialCharTok{==} \DecValTok{0}\NormalTok{),}
    \AttributeTok{sens =}\NormalTok{ tp }\SpecialCharTok{/}\NormalTok{ (tp }\SpecialCharTok{+}\NormalTok{ fn),}
    \AttributeTok{fpr  =}\NormalTok{ fp }\SpecialCharTok{/}\NormalTok{ (fp }\SpecialCharTok{+}\NormalTok{ tn)}
\NormalTok{  ) }\SpecialCharTok{|\textgreater{}}
  \FunctionTok{ungroup}\NormalTok{()}
\end{Highlighting}
\end{Shaded}

\begin{Shaded}
\begin{Highlighting}[]
\FunctionTok{ggplot}\NormalTok{(roc, }\FunctionTok{aes}\NormalTok{(fpr, sens)) }\SpecialCharTok{+}
  \FunctionTok{geom\_path}\NormalTok{(}\AttributeTok{linewidth =} \DecValTok{1}\NormalTok{) }\SpecialCharTok{+}
  \FunctionTok{geom\_abline}\NormalTok{(}\AttributeTok{linetype =} \DecValTok{2}\NormalTok{, }\AttributeTok{colour =} \StringTok{"grey50"}\NormalTok{) }\SpecialCharTok{+}
  \FunctionTok{coord\_equal}\NormalTok{() }\SpecialCharTok{+}
  \FunctionTok{labs}\NormalTok{(}\AttributeTok{x =} \StringTok{"False positive rate (1 {-} specificity)"}\NormalTok{,}
       \AttributeTok{y =} \StringTok{"Sensitivity"}\NormalTok{)}
\end{Highlighting}
\end{Shaded}

\subsection{5. Concluding statement}\label{concluding-statement}

\begin{quote}
With a cutoff of \texttt{cutoff}, the biomarker had sensitivity
\texttt{round(sens,\ 2)}, specificity \texttt{round(spec,\ 2)}, PPV
\texttt{round(ppv,\ 2)}, and NPV \texttt{round(npv,\ 2)}. The positive
likelihood ratio was \texttt{round(lrp,\ 2)} and the negative
\texttt{round(lrn,\ 2)}. A pre-test probability of 10\% becomes
\texttt{round(post\_prob\_pos,\ 2)} after a positive test and
\texttt{round(post\_prob\_neg,\ 3)} after a negative test.
\end{quote}

A single cutoff collapses a rich continuous score into two states. The
ROC curve shows the trade-off across all cutoffs; the area under it
summarises overall discrimination without committing to a threshold.

\subsection{Common pitfalls}\label{common-pitfalls}

\begin{itemize}
\tightlist
\item
  Quoting a single cutoff's sensitivity and specificity as if they were
  fixed properties of the test, ignoring that a different cutoff gives
  different numbers.
\item
  Confusing sensitivity with PPV in everyday speech.
\item
  Forgetting that PPV and NPV depend on prevalence.
\item
  Using an ROC to compare tests with different prevalence in each
  sample.
\end{itemize}

\subsection{Further reading}\label{further-reading}

\begin{itemize}
\tightlist
\item
  Altman DG \& Bland JM, \emph{Diagnostic tests} series, BMJ.
\item
  Pepe MS. \emph{The Statistical Evaluation of Medical Tests for
  Classification and Prediction}.
\end{itemize}

\subsection{Session info}\label{session-info}

\begin{Shaded}
\begin{Highlighting}[]

\end{Highlighting}
\end{Shaded}

\subsection{See also --- chapter index}\label{see-also-chapter-index}

\begin{itemize}
\tightlist
\item
  \href{../index.Rmd}{Methods in this chapter}
\item
  \href{../../../ch00-find-your-method.Rmd}{Find your method (Chapter
  0)}
\end{itemize}

\end{document}
